\documentclass[11pt]{article}
\usepackage{amsmath,amssymb,amsthm,hyperref}
\newtheorem{theorem}{Theorem}
\newtheorem{lemma}{Lemma}
\title{Positive Equality Modules and Universal Matching Covers}
\author{}
\date{November 2026}
\begin{document}
\maketitle

\paragraph{Status and scope.}
The finite results below are complete written proofs, not Lean proofs.
They give a different universal logarithmic-saving compiler on a
two-one-hot promised set. They do not improve the worst-case head
complexity on the whole Boolean cube. All intermediate denominators are
positive on the entire cube and have nonnegative slopes. Perturbing
zero slopes gives actual original-model heads with the same strict grid
signs and the same head count. The exact Boolean-axis identities belong
to the unperturbed intermediate modules. The missing global transfer is
identified explicitly below.

\begin{lemma}[A nonnegative equality module]
For $t>1$ and $X,Y>0$, put
\[
 K_t(X,Y)=\frac{tX}{tX+Y}-\frac{X}{X+Y}
 =\frac{(t-1)XY}{(tX+Y)(X+Y)}.
\]
For $\varepsilon>0$ and $Y\geq\varepsilon$, the axis-subtracted module
$G_t(X,Y;\varepsilon)=K_t(X,Y)-K_t(X,\varepsilon)$ satisfies
\[
 G_t(X,Y;\varepsilon)=
 \frac{(t-1)X(Y-\varepsilon)(tX^2-\varepsilon Y)}
 {(tX+Y)(X+Y)(tX+\varepsilon)(X+\varepsilon)}.
\]
The same formula holds at $X=0$, where all four denominators are
positive and both sides vanish.
\end{lemma}
\begin{proof}
Put the two fractions in the definition of $K_t$ over a common
denominator. For the second identity, subtract the resulting expressions.
The numerator remaining after factoring $(t-1)X$ is
\[
 Y(tX+\varepsilon)(X+\varepsilon)
 -\varepsilon(tX+Y)(X+Y)
 =(Y-\varepsilon)(tX^2-\varepsilon Y).
\]
Every denominator is nonzero under the hypotheses, so the algebra is
valid. At $X=0$ both original numerators vanish.
\end{proof}

Let $z\in\{0,1\}^r$, $w\in\{0,1\}^s$ and choose positive weights
$x_i,y_j$. Substitute
\[
 X=\sum_i x_i z_i,\qquad Y=\varepsilon+\sum_j y_jw_j.
\]
This gives four affine ratios with positive constant coefficients and
nonnegative slopes. It vanishes on both full Boolean axes.
If $x_{\min}=\min_i x_i$ and
$\varepsilon(\varepsilon+\sum_j y_j)<t x_{\min}^2$, it is nonnegative
on the entire cube: either $X=0$ or $X\geq x_{\min}$, and the
factorization proves the claim.

For $t=2$, $K_2(x,x)=1/6$. If $x/y\geq B$ or $y/x\geq B$, with
$B>1$, then $0<K_2(x,y)\leq1/B$, because
\[
 \frac{xy}{(2x+y)(x+y)}
 \leq \min\left\{\frac{x}{y},\frac{y}{2x}\right\}.
\]
Thus exponentially separated labels give a sharp equality test at
constant head cost, in the positive orientation.

\begin{theorem}[Universal matching array]
Let $M\geq64$, let $q$ be a positive integer with
$q\leq M/(5\log M)$, and put $H=2M$. There is an array
$a:[M]\times[H]\longrightarrow[q]$ such that, for every request vector
$u\in[q]^M$, distinct slots $h_1,\ldots,h_M$ can be chosen with
$a(g,h_g)=u_g$ for every $g$.
\end{theorem}
\begin{proof}
Choose all entries independently and uniformly. For a fixed request
vector, Hall's theorem fails only if some set $T$ of $t$ groups has
fewer than $t$ available slots. Equivalently, there is a set $U$ of
$H-t+1$ slots on which every group in $T$ misses its requested symbol.
To handle every request vector at once, it suffices to avoid pairs
$T,U$ for which every row in $T$ misses at least one alphabet symbol
across $U$. Independence between rows and a union bound over the
missing symbol give probability at most
\[
 \left[q(1-1/q)^{|U|}\right]^t
 \leq \left[q e^{-M/q}\right]^t.
\]
For $q=1$ a missing symbol is impossible, so the assertion is immediate.
For $q>1$, the displayed estimate uses
$|U|=2M-t+1\geq M+1$ and $1-x\leq e^{-x}$.
There are
$\binom Mt\binom{2M}{t-1}\leq(2M^2)^t$
such pairs. Hence the probability of any failure is at most
\[
 \sum_{t=1}^M [2M^2q e^{-M/q}]^t.
\]
The expression in brackets is at most
$2/(5M^2\log M)<1/2$, by the stated bound on $q$.
The finite geometric sum is therefore strictly less than one.
There exists an array with no failure. For such an array, Hall's
condition holds for every request vector, and Hall's theorem gives
the required matching.
\end{proof}

\begin{theorem}[Universal equality cover]
Take $q=2^k\leq M/(5\log M)$ and $r=Mk$. There are $H=2M$
fixed partitions of the $r$ rows such that every subset of rows is
the union of at most one block chosen from each partition.
Consequently every $r$ by $s$ Boolean matrix, for arbitrary finite
$s$, is a union of $H$ vertex-disjoint-biclique families.
\end{theorem}
\begin{proof}
Divide the rows into $M$ groups of size $k$, and identify the alphabet
$[q]$ with subsets of a group. In partition $h$, make the offered
subset $a(g,h)$ one block within group $g$, and its complement another
block, omitting empty blocks. Different groups are disjoint, so these
are partitions.

Given a target subset, let $u_g$ be its intersection with group $g$.
Use the universal matching array to choose distinct slots $h_g$.
At slot $h_g$ select the offered block $u_g$ if it is nonempty.
Select no block in the remaining slots. The selected blocks have
exactly the required union. Apply the same construction independently
to each matrix column. For a fixed slot, columns choosing the same
block are joined to all rows of that block. Since the row blocks
and column classes are disjoint, this is a vertex-disjoint union
of bicliques.
\end{proof}

\begin{theorem}[Positive-orientation grid compiler]
In the preceding setting, every strict sign table on the $r$ by $s$
two-one-hot grid has a representation by at most $4H+3$ positive-
orientation affine ratios and a bias, with exact zero on both full
Boolean axes before strict-slope perturbation. Its denominators are
positive on the entire $(r+s)$-cube. The strict grid signs have an
actual admissible-head realization with the same head count.
In particular the promised-set cost is $O(r/\log r)$.
\end{theorem}
\begin{proof}
Use the universal partitions. In slot $h$, assign label $g$ to the
offered block in group $g$, and label $M+g$ to its complement.
For column $j$, assign slot label $g$ when that slot is matched to
the nonempty requested subset in group $g$; assign label zero to
all unused slots. Thus a row has a matching label in some slot
if and only if its table entry is positive. It has exactly one
matching label when positive.

Choose $B>48H$, and set $x_{ih}=B^{\text{row label}}$,
$y_{jh}=B^{\text{slot label}}$. All these weights are at least one.
For each slot form the four-ratio module
$G_2(X_h,Y_h;\varepsilon)$ above.
Choose a common positive $\varepsilon$ so small that
\[
 \varepsilon<\frac1{96H},
 \qquad
 \varepsilon\left(\varepsilon+\sum_j y_{jh}\right)<2
 \quad\hbox{for every }h.
\]
This is possible because the slot and column sets are finite.
All four-ratio modules are then nonnegative throughout the cube.

On a grid point, replacing $y$ by $y+\varepsilon$ changes
$K_2(x,y)$ by at most $3\varepsilon/(2x)$: apply the elementary
identity $a/(a+y)-a/(a+y+\varepsilon)
=a\varepsilon/[(a+y)(a+y+\varepsilon)]$ to the two terms.
Also $K_2(x,\varepsilon)\leq\varepsilon/(2x)$.
Thus the repaired grid value differs from $K_2(x,y)$ by at most
$2\varepsilon$. The sum of all $H$ repaired modules differs by
less than $1/48$ from the unrepaired sum.

For a positive entry, the latter sum is at least $1/6$.
For a negative entry, it is less than $H/B<1/48$.
To implement a constant threshold without breaking axis zeros, put
$u=\sum_i z_i$, $v=\sum_jw_j$ and use
\[
 T(u,v)=1-\frac1{1+u}-\frac1{1+v}+\frac1{1+u+v}.
\]
It vanishes on both full axes and equals $1/3$ on the grid.
Subtract $T/4$. The final score is at least $7/48-1/12>0$ on
positive entries and less than $1/24-1/12<0$ on negative entries.
This uses $4H+3$ ratios and one bias.

Every denominator has positive constant coefficient and nonnegative
variable coefficients on the full cube. Add a common sufficiently
small positive multiple of $\sum_i z_i+\sum_jw_j$ to each denominator.
The coefficients become strictly positive, and finite-grid continuity
preserves the strict margins. The original
\texttt{FracAtom.ofPositiveAffineRatio} interface then gives actual
admissible heads. This last perturbation need not retain axis zeros;
the theorem distinguishes the two conclusions.

For sufficiently large $M$, take
$k=\lfloor\log_2(M/(5\log M))\rfloor$. Then $k=\Theta(\log M)$,
$r=Mk$, and $H=2M=O(r/\log r)$. Rounding or adding dummy rows gives
the same eventual order for arbitrary row counts.
\end{proof}

\paragraph{Exact remaining global obstruction.}
This supplies the positive orientation already known to be insufficient
for the established downward-star gluing. Substituting deletion
coordinates $z_i=1-x_i$ changes its slopes to negative ones, while the
outside-downset suppression penalty has positive slopes. The resulting
mixed slopes are not admitted by either original affine-ratio interface.
Nonnegativity of the equality modules does not repair this.

Indeed, adding a common penalty $L$ to the four denominators changes
the axis-subtracted expression to
$K_t(X,Y+L)-K_t(X,\varepsilon+L)$, whose sign factor is
\[
 (Y-\varepsilon)\bigl(tX^2-(\varepsilon+L)(Y+L)\bigr).
\]
For fixed $X>0$ and $Y>\varepsilon$, this is negative for all sufficiently
large $L$. Thus a claim that positive equality modules remain
nonnegative under ordinary suppression is false.
The known upward-star ordering has the complementary mixed-slope
problem. A model-preserving global assembly or a compact reverse
realization of the equality-cover mechanism has not been proved.

\paragraph{Literature transfer check.}
Erd\H{o}s and Pyber, \emph{Covering a graph by complete bipartite graphs},
Discrete Mathematics 170 (1997), 249--251,
\url{https://doi.org/10.1016/S0012-365X(96)00124-0},
bound the number of bicliques incident with each vertex by
$O(n/\log n)$. This is a local load bound, not a bound on the total
number of bicliques or on their simultaneous disjoint-union colors.
It cannot be substituted directly for the equality cover above.
The relevant cover notion is a union of vertex-disjoint biclique
families, called a bipartite $P_4$-free cover; see
\url{https://www.cs.purdue.edu/homes/hmaji/papers/BMMM20.pdf}.
The universal matching proof here establishes the needed finite cover
directly and does not assume that the local-load theorem supplies it.

\end{document}
